% ===== PRÉAMBULE CANONIQUE — à coller VERBATIM en tête de CHAQUE .tex =====
\documentclass[11pt,a4paper]{article}
\usepackage[utf8]{inputenc}\usepackage[T1]{fontenc}\usepackage{lmodern}
\usepackage[french]{babel}
\usepackage{amsmath,amssymb,mathtools}\usepackage{icomma}
\usepackage[version=4]{mhchem}          % \ce{...} : équations & formules
\usepackage{siunitx}\sisetup{locale=FR} % \qty{}{}, \si{}, \ang{}
\usepackage{chemfig}                    % \chemfig{...} : structures développées
\usepackage{xcolor}\usepackage{colortbl}
\usepackage{tikz}\usepackage{pgfplots}\pgfplotsset{compat=1.18}
\usepackage[most]{tcolorbox}\usepackage{enumitem}\usepackage{array,booktabs}
\usepackage[a4paper,margin=2.2cm]{geometry}
\usetikzlibrary{arrows.meta,calc,angles,quotes,positioning,patterns,backgrounds,shapes.geometric}
\definecolor{primary}{RGB}{222,49,99}\definecolor{primarydark}{RGB}{150,22,55}
\definecolor{defcol}{RGB}{0,114,178}\definecolor{methcol}{RGB}{52,92,148}
\definecolor{excol}{RGB}{0,135,90}\definecolor{attncol}{RGB}{200,40,55}
\tcbset{boxbase/.style={enhanced,breakable,boxrule=0.9pt,arc=2.5pt,
  left=3mm,right=3mm,top=2.3mm,bottom=2.3mm,fonttitle=\bfseries,coltitle=white}}
% --- boîtes TUTORIEL (noms EXACTS, ne pas renommer) ---
\newtcolorbox{cle}[1][]{boxbase,colback=defcol!6,colframe=defcol,
  title={\textsc{À retenir}\ifx\relax#1\relax\relax\else\textmd{ : \itshape #1}\fi}}
\newtcolorbox{methode}[1][]{boxbase,colback=methcol!7,colframe=methcol,sharp corners,
  title={\textsc{Méthode}\ifx\relax#1\relax\relax\else\textmd{ --- \itshape #1}\fi}}
\newtcolorbox{exemplebox}[1][]{boxbase,colback=excol!5,colframe=excol,
  title={\textsc{Exemple}\ifx\relax#1\relax\relax\else\textmd{ : \itshape #1}\fi}}
\newtcolorbox{piege}[1][]{boxbase,colback=attncol!6,colframe=attncol,
  title={\textsc{Piège}\ifx\relax#1\relax\relax\else\textmd{ : \itshape #1}\fi}}
\newtcolorbox{remarque}[1][]{boxbase,colback=black!4,colframe=black!45,
  title={\textsc{Remarque}\ifx\relax#1\relax\relax\else\textmd{ : \itshape #1}\fi}}
% --- boîtes EXERCICE / CORRIGÉ (noms EXACTS, auto-comptées) ---
\newtcolorbox[auto counter]{exo}[1][]{boxbase,colback=primary!4,colframe=primary,
  title={\textsc{Exercice}~\thetcbcounter\ifx\relax#1\relax\relax\else\textmd{ --- \itshape #1}\fi}}
\newtcolorbox[auto counter]{corrige}[1][]{boxbase,colback=excol!4,colframe=excol,
  title={\textsc{Corrigé}~\thetcbcounter\ifx\relax#1\relax\relax\else\textmd{ --- \itshape #1}\fi}}
\newcommand{\R}{\mathbb{R}}\newcommand{\N}{\mathbb{N}}\newcommand{\Z}{\mathbb{Z}}\newcommand{\ds}{\displaystyle}
% (un \usepackage{fancyhdr}+en-tête est possible ; il est ignoré côté web)
% =========================================================================
\begin{document}

\begin{center}{\large\bfseries\color{primarydark} Thème 5 — Corrigé Moyen}\end{center}

\begin{corrige}[préparer une solution diluée]
\begin{enumerate}[label=\alph*)]
  \item $V_1 = \dfrac{C_2 V_2}{C_1} = \dfrac{0{,}25\times 100}{5{,}0} = \SI{5.0}{\milli\litre}$.
  \item On prélève $\SI{5.0}{\milli\litre}$ de solution mère, on les verse dans une fiole jaugée de
        $\SI{100}{\milli\litre}$, puis on complète avec de l'eau jusqu'au trait de jauge.
\end{enumerate}
\end{corrige}

\begin{corrige}[dilution impossible]
La concentration visée ($\SI{4.0}{\mole\per\litre}$) est \emph{supérieure} à celle de la mère
($\SI{2.0}{\mole\per\litre}$). Or une dilution ne peut qu'\textbf{abaisser} la concentration
($C_2 < C_1$) : l'opération est \textbf{impossible}.
\end{corrige}

\begin{corrige}[fraction massique et pourcentage]
\begin{enumerate}[label=\alph*)]
  \item $w_{\ce{A}} = \dfrac{50}{200} = 0{,}25$ ; $w_{\ce{B}} = \dfrac{150}{200} = 0{,}75$ (somme $=1$).
  \item \ce{A} : $25\,\%$ en masse ; \ce{B} : $75\,\%$ en masse.
\end{enumerate}
\end{corrige}

\begin{corrige}[fraction molaire à partir des masses]
\begin{enumerate}[label=\alph*)]
  \item $n(\ce{H2O}) = \dfrac{90}{18} = \SI{5}{\mole}$ ; $n(\text{glucose}) = \dfrac{180}{180} = \SI{1}{\mole}$.
  \item $n_{\text{tot}} = \SI{6}{\mole}$ : $\chi_{\ce{H2O}} = \dfrac{5}{6} \approx 0{,}833$,
        $\chi_{\text{glucose}} = \dfrac{1}{6} \approx 0{,}167$. Somme : $0{,}833 + 0{,}167 = 1$.
\end{enumerate}
\end{corrige}

\begin{corrige}[molalité d'une solution de soude]
\begin{enumerate}[label=\alph*)]
  \item $n(\ce{NaOH}) = \dfrac{4{,}0}{40} = \SI{0.10}{\mole}$.
  \item $\SI{200}{\gram} = \SI{0.20}{\kilo\gram}$ : $m = \dfrac{0{,}10}{0{,}20} = \SI{0.50}{\mole\per\kilo\gram}$.
\end{enumerate}
\end{corrige}

\end{document}
